% !PRINTED 791
\weberpage{791}

\noindent
Normalformen der linearen Substitutionsgruppen II, 181.\\
Normalgleichung I, 465.\\
Normalkörper I, 464; II, 588.\\
Normaltheiler einer Gruppe I, 511; II, 11.\\
Null I, 16.

\medskip
\noindent
\textbf{O}bere Grenze für die Wurzeln I, 316.\\
Octa\"edergruppe II. 216.\\ % [sic] gedruckter Punkt nach der Bandziffer
--- im Congruenzkörper II, 277.\\
Orthogonale Substitution II, 184.

\medskip
\noindent
\textbf{P}artialdiscriminante II, 597.\\
Partialgrundideal II, 585, 597.\\
Partialnorm II, 584.\\
Partialspur II, 585.\\
Partialresolventen I, 512.\\
\gesperrt{Pell}'sche Gleichung I, 395.\\
Periode einer Permutation I, 500.\\
Perioden der Kreistheilung I, 557, 579; II, 74.\\
--- der Wurzeln einer cyklischen Gleichung I, 545.\\
--- der reducirten Zahlen I, 388.\\
Periodische Kettenbrüche I, 387.\\
Permutationen I, 64, 473, 500.\\
--- ihre analytische Darstellung I, 613; II, 299.\\
--- erster und zweiter Art I, 65, 496.\\
Permutationsgruppen I, 475, 489; II, 117.\\
--- $168^{\mathrm{sten}}$ Grades von sieben Ziffern II, 473.\\
Polaren I, 188.\\
Pole linearer gebrochener Substitutionen II, 196.\\
--- einer Gruppe II, 198.\\
--- --- --- im Congruenzkörper II, 274.\\
--- einer ternären Gruppe II, 438.\\
Poly\"edergruppen II, 202.\\
--- der zweiten Art II, 234.\\
--- im Congruenzkörper II, 273.\\
Polynomialco\"efficienten I, 51.\\
Polynomischer Lehrsatz I, 50.\\
Positive Einheiten im Kreistheilungskörper II, 742.\\
Potenzreste I, 430.\\
Potenzsummen I, 140.

Primäre Theiler eines Kreistheilungskörpers II, 71.\\
--- --- einer \gesperrt{Abel}'schen Gruppe II, 72.\\
Primfactoren der Zahlen eines algebraischen Körpers II, 523.\\
--- der natürlichen Primzahlen II, 576.\\
--- der Kreistheilungsresolventen II, 635.\\
Primfunctionale II, 515, 576.\\
--- relative II, 514.\\
Primideale II, 550.\\
--- im quadratischen Körper II, 604.\\
--- im relativ normalen Körper II, 588.\\
--- in den Theilern eines Normalkörpers II, 599.\\
--- in den Kreistheilungskörpern II, 619, 622.\\
--- im reellen Kreistheilungskörper II, 643.\\
Primitive Congruenzwurzeln I, 428; II, 546.\\
--- Einheitswurzeln I, 410.\\
--- Functionen I, 25; II, 497.\\
--- Wurzeln von Primzahlen I, 429.\\
--- --- von Primzahlquadraten II, 57.\\
--- --- von Primfunctionalen II, 546.\\
--- --- von Primidealen II, 591.\\
--- --- eines Congruenzkörpers II, 248.\\
--- und imprimitive Charaktere in Kreistheilungskörpern II, 720.\\
--- --- --- Gruppen I, 484.\\
--- --- --- Körper I, 463.\\
--- --- --- Formenprobleme II, 180.\\
Primzahlen I, 1.\\
--- relative I, 359.\\
--- im Körper $R(i)$ I, 587.\\
--- im Körper der dritten Einheitswurzeln I, 593.

\medskip
\noindent
\textbf{Q}uadrate im Congruenzkörper II, 248.\\
Quadratische Formen I, 179.\\
--- Gleichungen I, 116.\\
--- Irrationalzahlen I, 377.\\
--- Reste I, 444.\\
--- Körper II, 601.\\
Quadratur des Kreises II, 760.

\medskip
\noindent
\textbf{R}ationale Wurzeln einer Gleichung I, 403.
